\section{Introducción y problema de investigación}
La instrumentación didáctica organiza contenidos, actividades, evaluación y tiempos en la educación superior tecnológica. El PID original administraba entregas documentales de instrumentaciones y reportes. La última versión recibía instrumentaciones XLSX, aunque algunos borradores de la sesión 1 describían incorrectamente PDF como formato de entrada. El sistema también contempla reportes de salida PDF.

El problema de investigación consiste en determinar cómo evolucionar del registro documental al seguimiento estructurado por unidad temática y actividad, con fechas programadas y reales, solicitudes de cambio autorizadas y conservación de versiones.

\subsection{Pregunta y objetivo}
¿Cómo diseñar un modelo de mejora funcional del PID que integre seguimiento programático y gestión autorizada de modificaciones a la planeación docente, fortaleciendo control y trazabilidad? El objetivo general es analizar y diseñar un modelo de mejora funcional con seguimiento por unidad y actividad, calendarios y modificaciones autorizadas.

\subsection{Alcance y continuidad de la sesión 1}
El objeto de estudio comprende planeación, seguimiento y cambios; excluye control escolar, calificaciones individuales y finanzas. La revisión crítica de antecedentes, análisis comparativo de ciclo de vida, factibilidad y referencias completas de P1 están pendientes de integración fiel desde su versión definitiva; este apartado es un puente editorial, no reemplaza P1.
